%% ======================================================================
\chapter{PAC-Bayesian learning of ensembles}
\label{chap:algos}
%% ======================================================================

\section{Self-bounding learning}

A PAC-Bayesian bound holds \emph{simultaneously for all posteriors}. We can
therefore use it as a learning objective:
\[
Q^\star\in\argmin_{Q\in\M(\Hcal)}\ \mathcal B\big(\widehat{\text{stat}}_S(Q),\KL(Q\|P),m,\delta\big),
\]
where $\widehat{\text{stat}}_S(Q)$ is an empirical quantity (Gibbs risk, joint
error, disagreement\dots). The resulting predictor comes with its own
generalization guarantee: the value of the bound at $Q^\star$ is a valid
\emph{risk certificate}, without any test set. Such algorithms are called
\emph{self-bounding} \citep{freund1998self}; in the PAC-Bayesian context, see
\citet{germain2015risk} and \citet{viallard2021self}. In this chapter we apply this principle to the
ensemble methods of the course.

\begin{remark}
A bound minimizer is not necessarily a risk minimizer: the bound is a
worst-case, high-probability statement. Minimizing it is a principled form of
regularization, with the KL as regularizer and $P$ as the reference point.
The quality of the result depends on how well the bound reflects the actual
behaviour of the predictor --- which, for majority votes, depends on the
\emph{order} of the bound (first vs.\ second order).
\end{remark}

\section{Bagging and random forests: out-of-bag PAC-Bayes}\label{sec:oob}

\subsection{The split construction}

The prior must not depend on the data used to evaluate the bound. With
bagging this looks problematic, since every tree is learned from $S$. But
each tree $h_t$ is learned on a bootstrap sample $S_t$ only, and the
\emph{out-of-bag} (OOB) examples $T_t=S\setminus S_t$ were not used to build it:
from the point of view of $h_t$, $T_t$ is a fresh test sample
(\cref{fig:oob}).

\begin{figure}[t]
\centering
\begin{tikzpicture}[scale=0.55]
\foreach \t/\pat in {1/{1,1,0,1,0,1,1,0,1,1,0,1,1,1,0,1},2/{0,1,1,1,1,0,1,1,0,1,1,0,1,1,1,0},3/{1,0,1,1,1,1,0,1,1,0,1,1,0,1,1,1}}{
  \node[left] at (0,-\t*1.2+0.35) {\small $h_{\t}$};
  \foreach \b [count=\i] in \pat {
    \ifnum\b=1
      \fill[pbblue!45] (\i-1,-\t*1.2) rectangle (\i,-\t*1.2+0.7);
    \else
      \fill[pborange!70] (\i-1,-\t*1.2) rectangle (\i,-\t*1.2+0.7);
    \fi
    \draw[white] (\i-1,-\t*1.2) rectangle (\i,-\t*1.2+0.7);
  }
}
\node[left] at (0,0.35) {\small $S$};
\foreach \i in {1,...,16}{\draw (\i-1,0) rectangle (\i,0.7); \node[font=\tiny] at (\i-0.5,0.35) {$z_{\i}$};}
\fill[pbblue!45] (17.5,-1.1) rectangle (18.3,-0.6); \node[right,font=\small] at (18.4,-0.85) {in the bootstrap sample $S_t$ (training)};
\fill[pborange!70] (17.5,-2.1) rectangle (18.3,-1.6); \node[right,font=\small] at (18.4,-1.85) {out-of-bag $T_t=S\setminus S_t$ (evaluation)};
\node[right,font=\small,align=left] at (17.5,-3.3) {tandem loss of $(h_1,h_2)$: evaluated on $T_1\cap T_2$};
\end{tikzpicture}
\caption{Out-of-bag construction. Each tree is trained on its bootstrap
sample (blue) and its loss is estimated on its OOB sample (orange), which is
independent of the tree. The joint error of two trees is estimated on the
examples that are out-of-bag for both.}
\label{fig:oob}
\end{figure}

\begin{theorem}{PAC-Bayes with a split construction \citep{thiemann2017strongly,lorenzen2019pac}}{oob}
Let $h_1,\dots,h_n$ be voters, $h_t$ being learned on $S_t\subset S$ by any
procedure, and let $T_t=S\setminus S_t$ with $|T_t|\geq n_0$ for all $t$. Let
$\widehat L_t=\frac{1}{|T_t|}\sum_{(x,y)\in T_t}\ind{h_t(x)\neq y}$ be the OOB
loss of $h_t$ and $\pi$ a distribution on $\{1,\dots,n\}$ fixed in advance
(e.g.\ uniform). Then, w.p.\ $\geq1-\delta$ over $S\sim\D^m$ (and the
randomness of the bootstrap), for all $\rho\in\M(\{1,\dots,n\})$,
\begin{equation}\label{eq:oob}
\kl\Big(\sum_t\rho_t\widehat L_t\ \Big\|\ \sum_t\rho_t\RD(h_t)\Big)\leq\frac{\KL(\rho\|\pi)+\ln\frac{2\sqrt{n_0}}{\delta}}{n_0}.
\end{equation}
The same holds for the tandem losses $\widehat L_{tt'}$ computed on $T_t\cap T_{t'}$,
with $\KL(\rho^2\|\pi^2)=2\KL(\rho\|\pi)$ and $n_0\leq\min_{t,t'}|T_t\cap T_{t'}|$.
\end{theorem}

\begin{proofsketch}
We follow the proof of Theorem~\ref{thm:general} with ``hypotheses'' being the
\emph{indices} $t$, on which $\pi$ is a legitimate data-independent prior. Only
Step~3 changes: for each $t$, conditionally on $S_t$ the voter $h_t$ is fixed
and $T_t$ is an i.i.d.\ sample independent of it, so Maurer's lemma gives
$\E\,e^{|T_t|\kl(\widehat L_t\|\RD(h_t))}\leq2\sqrt{|T_t|}$. Since $|T_t|\geq n_0$,
Jensen's inequality yields $\E\,e^{n_0\kl(\widehat L_t\|\RD(h_t))}\leq(2\sqrt{|T_t|})^{n_0/|T_t|}\leq2\sqrt{n_0}$.
\end{proofsketch}

In practice, with standard bootstrap each OOB set contains about
$(1-1/m)^m\approx e^{-1}\approx37\%$ of the sample, and the intersection of two OOB
sets about $e^{-2}\approx14\%$. Drawing bootstrap samples of size $m/2$ instead
of $m$ increases these proportions to about $e^{-1/2}\approx61\%$ and
$e^{-1}\approx37\%$, at the price of slightly weaker trees (this is the choice
made in the figures and notebooks of this course).

\subsection{Minimizing the first-order bound}\label{sec:algo_fo}

With the PAC-Bayes-$\lambda$ form \eqref{eq:lambda} of the bound, for a fixed
$\lambda$ the optimal posterior is the Gibbs posterior on the OOB losses
(Proposition~\ref{prop:gibbs_post}), and for a fixed posterior the optimal $\lambda$ is
given by \eqref{eq:lambda_star}. This suggests an alternating minimization.

\begin{algorithm}[h]
\caption{PAC-Bayes-$\lambda$ minimization for a forest \citep{thiemann2017strongly}}\label{algo:fo}
\KwIn{OOB losses $\widehat L_1,\dots,\widehat L_n$, prior $\pi$, $n_0$, $\delta$}
$\lambda\leftarrow1$\;
\Repeat{convergence of $\lambda$}{
  $\rho_t\leftarrow\dfrac{\pi_t\,e^{-\lambda n_0\widehat L_t}}{\sum_{s}\pi_s\,e^{-\lambda n_0\widehat L_s}}$ for $t=1,\dots,n$\;
  $\lambda\leftarrow\dfrac{2}{\sqrt{2n_0\,\rho^\top\widehat L/\big(\KL(\rho\|\pi)+\ln\frac{2\sqrt{n_0}}{\delta}\big)+1}+1}$\;
}
\KwOut{$\rho$ and the certificate $2\,\klup\big(\rho^\top\widehat L,\frac{\KL(\rho\|\pi)+\ln(2\sqrt{n_0}/\delta)}{n_0}\big)$ on $\RD(B_\rho)$}
\end{algorithm}

\citet{thiemann2017strongly} proved that, under mild conditions, the
PAC-Bayes-$\lambda$ bound is strongly quasiconvex and that this procedure
converges to its global minimum. It produces excellent bounds on the
\emph{Gibbs} risk. However, as observed by \citet{lorenzen2019pac} and
\citet{masegosa2020second}, the resulting majority vote is often \emph{worse}
than the uniform vote: the Gibbs posterior concentrates on the few trees with
the smallest OOB loss, which are highly correlated, and the benefit of the vote
disappears (\cref{fig:rf_opt}). This is a direct consequence of the first-order
nature of the bound, which ignores the correlations between voters.

\subsection{Minimizing the tandem bound}

The tandem bound \eqref{eq:tnd_pb} with the PAC-Bayes-$\lambda$ relaxation
gives the objective
\begin{equation}\label{eq:tnd_objective}
\mathcal B_{\text{TND}}(\rho,\lambda)=4\left[\frac{\rho^\top\widehat{\mathbf L}\rho}{1-\frac\lambda2}+\frac{2\KL(\rho\|\pi)+\ln\frac{2\sqrt{n_0}}{\delta}}{\lambda\big(1-\frac\lambda2\big)n_0}\right],
\qquad\widehat{\mathbf L}=(\widehat L_{tt'})_{t,t'},
\end{equation}
which is quadratic in $\rho$ (plus the convex KL term). Its minimization has no
closed form but is easily performed by gradient descent on the softmax
parametrization $\rho=\mathrm{softmax}(\theta)$, alternated with the closed-form
update of $\lambda$ \citep{masegosa2020second}. Because the objective penalizes
pairs of voters that err \emph{together}, the optimal posterior keeps the
ensemble diverse: in \cref{fig:rf_opt} the tandem-optimized forest is as
accurate as the uniform one, and its certificate is tighter.

\begin{figure}[t]
\centering
\includegraphics[width=\linewidth]{rf_optimisation}
\caption{Forests of 100 trees (bootstrap samples of size $m/2$), on four
datasets. For the uniform posterior, the posterior minimizing the first-order
(FO) bound, and the posterior minimizing the tandem (TND) bound, we report the
test risk of the majority vote and the two OOB certificates ($\delta=0.05$).
Minimizing the FO bound improves the FO certificate but can seriously degrade
the vote (e.g.\ \texttt{digits}); minimizing the TND bound preserves diversity.
Generated with \texttt{pbtools.py} (notebook 4).}
\label{fig:rf_opt}
\end{figure}

\section{Boosting-like algorithms: learning the weights of a vote}

\subsection{AdaBoost from the PAC-Bayesian viewpoint}

AdaBoost greedily minimizes the exponential loss $\frac1m\sum_ie^{-y_if(x_i)}$
of $f=\sum_t\alpha_th_t$, an upper bound on the training error of the vote
\citep{freund1997decision}. In PAC-Bayesian terms, it learns a posterior
$Q\propto(\alpha_t)_t$ whose \emph{margins} $M_Q$ are large on $S$ (the
normalized margins are $y_if(x_i)/\|\alpha\|_1$), which is the explanation of
\citet{schapire1998boosting}. But nothing controls $\KL(Q\|P)$ or the
disagreement explicitly. The following algorithms instead learn the weights of
a set of (possibly boosted) voters by minimizing a C-bound.

\subsection{MinCq: minimizing the C-bound}

Recall that $\mathcal C_Q=1-\frac{(\E M_Q)^2}{\E M_Q^2}$. For a fixed value of the
first moment $\E M_Q=\mu$, minimizing the C-bound amounts to minimizing the
second moment of the margin. \citet{laviolette2011from} and
\citet{germain2015risk} proposed
\begin{equation}\label{eq:mincq}
\textbf{MinCq:}\qquad \min_{Q}\ \widehat{\E}_S\big[M_Q^2\big]\quad\text{s.t.}\quad\widehat{\E}_S\big[M_Q\big]=\mu,\quad Q\text{ quasi-uniform},
\end{equation}
where $\Hcal$ is \emph{self-complemented} ($h\in\Hcal\Rightarrow-h\in\Hcal$) and a
\emph{quasi-uniform} posterior satisfies $Q(h)+Q(-h)=\frac1{|\Hcal|/2}$. Since
$M_Q$ is linear in $Q$, $\widehat\E_S[M_Q^2]$ is a quadratic form in the weights
and \eqref{eq:mincq} is a \emph{quadratic program}. Two remarkable facts:
\begin{itemize}
\item on quasi-uniform posteriors, PAC-Bayesian bounds hold \emph{without the KL
      term} \citep{germain2015risk}, so that the empirical C-bound itself is controlled;
      the restriction is not a loss of expressiveness because any vote can be
      represented by a quasi-uniform posterior with the same predictions;
\item the hyperparameter $\mu$ (the targeted first moment of the margin)
      plays the role of a regularization parameter, selected by cross-validation.
\end{itemize}
MinCq has been extended to incorporate prior knowledge on the weights
\citep[P-MinCq,][]{bellet2014learning} and to a column-generation (boosting-like)
algorithm that adds voters one at a time \citep[CqBoost,][]{roy2016column}.

\subsection{Direct minimization of PAC-Bayesian C-bounds}

\citet{viallard2021self} proposed to minimize by (stochastic) gradient descent
the \emph{PAC-Bayesian} C-bounds themselves, i.e.\ the right-hand side of
\eqref{eq:cbound_pb} and tighter joint versions, over the posterior
$\rho=\mathrm{softmax}(\theta)$. The learned vote comes with a tight certificate
and competes with the state of the art. \citet{zantedeschi2021learning} go one
step further: they consider \emph{stochastic majority votes} whose weights are
themselves drawn from a Dirichlet posterior; for such votes the expected risk
can be computed (or tightly approximated) and bounded directly, avoiding the
loose Gibbs-to-vote relations of \cref{chap:mv}.

\section{Linear classifiers and kernel methods}

\subsection{Gaussian posteriors on linear classifiers}

Consider linear voters $h_{\mathbf v}(x)=\sign(\langle\mathbf v,x\rangle)$,
$\mathbf v\in\R^d$, a prior $P=\mathcal N(\mathbf 0,I_d)$ and posteriors
$Q_{\mathbf w}=\mathcal N(\mathbf w,I_d)$.
\begin{itemize}
\item \textbf{Majority vote.} By symmetry of the Gaussian around $\mathbf w$,
      $B_{Q_{\mathbf w}}(x)=\sign\big(\E_{\mathbf v}\sign\langle\mathbf v,x\rangle\big)=\sign\langle\mathbf w,x\rangle$:
      the vote of an infinite number of linear classifiers is the linear
      classifier $\mathbf w$.
\item \textbf{Gibbs risk.} For $\mathbf v\sim Q_{\mathbf w}$,
      $y\langle\mathbf v,x\rangle\sim\mathcal N\big(y\langle\mathbf w,x\rangle,\|x\|^2\big)$, hence
      \begin{equation}\label{eq:probit}
      \RS(G_{Q_{\mathbf w}})=\frac1m\sum_{i=1}^m\Phi\Big(-\frac{y_i\langle\mathbf w,x_i\rangle}{\|x_i\|}\Big),
      \end{equation}
      where $\Phi$ is the standard Gaussian CDF: the 0-1 loss is replaced by
      a smooth \emph{probit} loss of the normalized margin.
\item \textbf{Complexity.} $\KL(Q_{\mathbf w}\|P)=\frac12\|\mathbf w\|^2$.
\end{itemize}
Plugging these quantities into Catoni's bound \eqref{eq:catoni} and removing
the monotone transformation, one obtains the objective of the algorithm
PBGD \citep{germain2009pac}:
\begin{equation}\label{eq:pbgd}
\min_{\mathbf w\in\R^d}\ C\sum_{i=1}^m\Phi\Big(-\frac{y_i\langle\mathbf w,x_i\rangle}{\|x_i\|}\Big)+\frac12\|\mathbf w\|^2,
\end{equation}
to be compared with the soft-margin SVM $\min_{\mathbf w}C\sum_i\max(0,1-y_i\langle\mathbf w,x_i\rangle)+\frac12\|\mathbf w\|^2$:
PAC-Bayes \emph{explains} the SVM objective (regularization $=$ KL to the
prior, loss $=$ a surrogate of the Gibbs risk), and the norm of $\mathbf w$
plays the role of the inverse of a margin \citep{langford2002pac,mcallester2003simplified}.
\Cref{fig:linear} illustrates the trade-off: scaling $\mathbf w$ up makes the
Gibbs classifier behave like the vote (smaller Gibbs risk) but increases the KL.
The objective \eqref{eq:pbgd} is not convex; it is optimized by gradient
descent with several restarts, and has a kernelized version
\citep{germain2009pac}. Data-dependent Gaussian priors centred on an SVM
learned on part of the data give some of the tightest bounds for SVMs
\citep{ambroladze2006tighter,parrado2012pac}.

\begin{figure}[t]
\centering
\includegraphics[width=0.85\linewidth]{linear_gaussian}
\caption{Gaussian posterior $\mathcal N(\mathbf w,I)$ over linear classifiers:
40 sampled hyperplanes (grey) and the majority vote $\sign\langle\mathbf w,x\rangle$
(black). Multiplying $\mathbf w$ by $5$ does not change the vote but makes the
sampled voters agree with it: the Gibbs risk decreases while the KL grows
quadratically.}
\label{fig:linear}
\end{figure}

\subsection{Kernels as priors: random Fourier features}

For a shift-invariant kernel $k(x,x')=k(x-x')$, Bochner's theorem gives
$k(x-x')=\E_{\omega\sim p}\cos\big(\omega^\top(x-x')\big)$ where $p$ is the Fourier
transform of $k$ \citep{rahimi2007random}. \citet{letarte2019pseudo} interpret
this as a \emph{majority vote}: the voters are the functions
$h_\omega(x,x')=\cos(\omega^\top(x-x'))$ and the kernel is the vote under the
\emph{prior} $p$. Learning a kernel adapted to the task then amounts to learning a
\emph{posterior} $Q$ over frequencies; with $N$ frequencies $\omega_1,\dots,\omega_N$
drawn from $p$, the pseudo-posterior $Q(\omega_j)\propto e^{-\beta\,\widehat L_S(h_{\omega_j})}$
is exactly a Gibbs posterior, and the PAC-Bayesian bound certifies the
learned kernel. This is closely related to the boosting-with-random-Fourier-features
algorithm studied in the first project of the course.

\section{Neural networks and deep ensembles}

The same machinery applies to neural networks with a Gaussian posterior on
the weights (a \emph{stochastic} neural network). \citet{dziugaite2017computing}
optimized a PAC-Bayesian bound over such posteriors and obtained the first
\emph{non-vacuous} bounds for over-parametrized networks on MNIST;
\citet{perez2021tighter} showed that \emph{training} with PAC-Bayesian objectives
and data-dependent priors yields both accurate networks and tight certificates,
and \citet{lotfi2022pac} obtained non-vacuous bounds on large models with
compression-based priors. \emph{Deep ensembles} \citep{lakshminarayanan2017simple}
are majority votes of independently trained networks, to which the
second-order bounds of \cref{chap:mv} directly apply.

\section{Practical guidelines}

\begin{enumerate}
\item \textbf{Decide what is certified.} Gibbs classifier, majority vote, or a
      single deterministic model? Use the corresponding bound (\cref{chap:mv}).
\item \textbf{Fix the prior before looking at the evaluation data}: uniform over
      voters built on other data (split, OOB), Gaussian around a model learned on a
      separate sample, etc.
\item \textbf{Account for every choice made with the data}: a grid of $k$ values for
      $\lambda$, $C$ or $\sigma$ costs $\ln k$ via a union bound.
\item \textbf{Use the kl form} ($\klup$) to compute the final certificate, even if a
      relaxation was used to learn $Q$.
\item \textbf{Prefer second-order bounds} (tandem, C-bound, CCTND) for majority
      votes of diverse voters; they need enough data to pay the doubled KL.
\item \textbf{Report} the bound together with its ingredients (empirical term, KL,
      $m$, $\delta$) and compare it with the test error: a bound is a guarantee,
      not an estimate.
\end{enumerate}
